\chapter{Atlas of the literature}\label{app:atlas}
% STATUS: written 2026-11-06 (counts computed from bib/encyclopaedia.bib, the
% collected bibliography of this edition, and from the kind census in
% ledger/encyclopaedia-kinds.tsv. No external claim is made here; every number is
% reproducible by re-running the census, so the appendix moves with each
% edition.)
%
% The counts below are generated from the book's collected bibliography and
% refreshed with each edition. They are described in words, with no internal
% file name printed (owner, 2026-11-06).

The bibliography printed with this edition holds works from 1684 to this year.
This appendix describes what that collection looks like, so that a reader can
see the shape of the field at a glance.
The counts cover the works the text cites.
They do not measure the field itself, which is larger.
They measure the part of it this edition has read.

\section{Works by decade}

The collection is thin before 1900 and grows through the twentieth century.
Table~\ref{tab:atlas:decades} gives the exact count, and
Figure~\ref{fig:atlas:decades} draws it.

\hbfigure{backmatter/works-by-decade}{Works per decade. The single 1680s entry is
Leibniz, where the product rule is stated; the 2010s is the peak at 114.}{fig:atlas:decades}

Three things in the figure are worth naming. The first bar is alone for two
centuries: Leibniz is the only work this edition cites from the 1600s, and
nothing follows until the 1880s. The figure then rises unevenly, doubling about
every two decades after 1950. And the last two bars are nearly equal, which is
what a decade still in progress looks like.

\begin{table}[htbp]
  \centering
  \caption{Works in the bibliography, by decade of publication.}
  \label{tab:atlas:decades}
  \begin{tabular}{lr@{\qquad}lr}
    \toprule
    Decade & Works & Decade & Works \\
    \midrule
    1680s & 1   & 1960s & 13  \\
    1880s & 1   & 1970s & 27  \\
    1890s & 1   & 1980s & 26  \\
    1910s & 2   & 1990s & 50  \\
    1920s & 3   & 2000s & 59  \\
    1930s & 5   & 2011s & 114 \\
    1940s & 6   & 2020s & 105 \\
    1950s & 7   &       &     \\
    \bottomrule
  \end{tabular}
\end{table}

Three features stand out.
The first is the single entry in the 1680s, the work of Leibniz where the
product rule is stated.
The second is the slow start: fewer than ten works a decade before 1950.
The third is the recent weight.
Of the 456 works, 304 appeared after 2000, and the 2010s alone hold more
works than every decade before 1950 together.

The 2020s count is close behind the 2010s, and the decade is not over.

\section{Works by language}

Table~\ref{tab:atlas:language} gives the language of the collection, and
Figure~\ref{fig:atlas:language} shows when the collection changed language.

\hbfigure{backmatter/language-shift}{Works per decade by language. The gold
stratum is German and Latin, which dominates before 1940; the blue is English,
which is nearly all of the collection after it.}{fig:atlas:language}

The two strata barely overlap. German and Latin carry the subject to 1940, with
one late German work in the 1970s; English takes over in the 1930s and is the
whole collection from 1990. The two French works and the single Latin one sit
at the ends of that movement rather than outside it.

\begin{table}[htbp]
  \centering
  \caption{Works in the bibliography, by language.}
  \label{tab:atlas:language}
  \begin{tabular}{lr}
    \toprule
    Language & Works \\
    \midrule
    English & 408 \\
    German  & 9   \\
    French  & 2   \\
    Latin   & 1   \\
    \bottomrule
  \end{tabular}
\end{table}

Most of the collection is in English.
The German and Latin works are the older ones, from the period when the subject
was written in those languages, and the two French works are cited for the same
reason.
That distribution follows the history in Chapter~\ref{ch:origins}, where the
thread passes from Latin to German and then to English as the centres of the
subject move.

\section{Kinds of derivation-like map}

The kind census holds 200 rows. Twenty-seven of them are tokenizer
fragments rather than names, and they are listed with their reason in the
kind ledger of the book, so the book names 173 kinds of derivation-like map.
A number in that table or in the index of kinds is the number of works in
this atlas whose title names the kind. A title that names three kinds counts
once under each of the three, so the numbers do not add up to a number of works.
Table~\ref{tab:atlas:kinds} gives the eight that the literature treats most.

\begin{table}[htbp]
  \centering
  \caption{The eight most studied kinds, with the number of works in the census
  that name each.}
  \label{tab:atlas:kinds}
  \begin{tabular}{lr}
    \toprule
    Kind & Works \\
    \midrule
    generalized derivation & 592 \\
    skew derivation        & 179 \\
    higher derivation      & 145 \\
    Jordan derivation      & 135 \\
    Lie derivation         & 107 \\
    triple derivation      & 98  \\
    nilpotent derivation   & 85  \\
    $*$-derivation         & 82  \\
    \bottomrule
  \end{tabular}
\end{table}

The ordering is not a ranking of importance.
It reflects how the literature itself distributes its attention, and it
explains the length of the chapters in this book.
Generalized derivations, which include the skew and the twisted cases, occupy
Chapter~\ref{ch:generalized-twisted} for that reason.
Local and 2-local derivations, which ask when a map that is a derivation at each
point is a derivation everywhere, occupy
Chapter~\ref{ch:local-characterizations}.

\section{Kinds by chapter}

Table~\ref{tab:atlas:kinds-chapters} groups the same census by the chapter that
treats each kind.

\begin{table}[htbp]
  \centering
  \caption{Works in the kind census, grouped by the chapter that treats the
  kind.}
  \label{tab:atlas:kinds-chapters}
  \begin{tabular}{lr@{\qquad}lr}
    \toprule
    Chapter & Works & Chapter & Works \\
    \midrule
    Generalized and twisted & 993 & Basic theory          & 30 \\
    Prime and semiprime     & 306 & Polynomial and Weyl   & 24 \\
    Banach and operator     & 301 & Matrix and triangular & 13 \\
    Lie-type                & 230 & Functional identities & 13 \\
    Other systems           & 184 & Unclassified          & 12 \\
    Higher and multi        & 181 & Cohomology            & 5  \\
    Local maps              & 176 & Origins               & 3  \\
    Commutative             & 163 & Group and Hopf        & 2  \\
    Jordan-type             & 144 &                       &    \\
    Nonassociative          & 136 &                       &    \\
    Lie algebras            & 66  &                       &    \\
    Applications            & 37  &                       &    \\
    \bottomrule
  \end{tabular}
\end{table}

The two largest groups are the generalized and the prime-ring kinds.
Together they hold 1,299 of the 3,019 named occurrences in the census, which is
43 per cent of it.
A work can name more than one kind, so the census counts occurrences rather
than works, and its total is larger than the number of works in the
bibliography.
The smallest are the group and Hopf kinds, with 2 works, which are the most
recent arrivals in the literature and which Chapter~\ref{ch:group-hopf} treats
as an open frontier.

\section{How the collection has grown}

Figure~\ref{fig:atlas:cumulative} draws the same dates cumulatively, which is
the clearer way to see the shape.

\hbfigure{backmatter/cumulative-growth}{Works in the collection up to each
decade. The curve is flat for two centuries, then steep; the halfway mark falls
at the end of the 2000s.}{fig:atlas:cumulative}

The curve has an elbow, and the elbow is the point of the section. Half of the
bibliography was published after the median year, so this edition has read
mostly recent work. That is a choice the subject forces: the questions collected in
Section~\ref{sec:basic-theory:open-problems-by-class} are being worked on now, and a
encyclopaedia that described only settled work would miss them. The cost is visible in the same
curve, which barely leaves the axis before 1950.

The dates in the bibliography describe a subject that grew slowly and then
quickly.
Table~\ref{tab:atlas:growth} gives four measures of that growth.

\begin{table}[htbp]
  \centering
  \caption{Four measures of the collection, computed from the dates in the
  bibliography.}
  \label{tab:atlas:growth}
  \begin{tabular}{lr}
    \toprule
    Measure & Value \\
    \midrule
    Works with a date            & 456 of 456 \\
    Earliest and latest          & 1684 and 2026 \\
    Median year                  & 2010 \\
    Mean year                    & 1998 \\
    Works published after 2000   & 304 \\
    Works published before 1950  & 19 \\
    Works published since 2010   & 219 (half the collection) \\
    \bottomrule
  \end{tabular}
\end{table}

The mean lies eight years below the median, so the distribution has a long
early tail: a few very old works pull the mean back without moving the middle.
The median year of 2010 means that half the collection was published in the
last twenty years.
The peak decade is the 2010s with 114 works, and the three most productive
decades hold 278 of the 420 works, which is 66 per cent of them.

Every entry carries a year, so the table describes the whole collection and
no correction is needed for undated works.

\section{How the atlas grows}

The collection is not fixed.
Each edition adds the works read since the last one, and the counts above are
counted again when it does.
A work enters the bibliography only when its text has been read from the page,
so the atlas records the reading, not the field.

The counts are produced by one program from two records, the bibliography and
the kind census, so a new download that adds a work changes the tables at the
next rebuild without any hand editing.
That is what makes the appendix a living part of the edition rather than a
snapshot taken once.

The gap between the collection and the field is the work of later editions, and
Appendix~\ref{app:questions} lists the questions that gap leaves open.
